%==============================================================================
% Section 6: Four-topology admissibility classification
%==============================================================================
\section{Four-topology admissibility classification}
\label{sec:classification}

The local admissibility classification that integrates the Hamiltonian
admissibility of Section~\ref{sec:hamiltonian} and the P-channel diagnostics
of Section~\ref{sec:pchannel} is collected in
Table~\ref{tab:four_topology_admissibility}.

\begin{table}[H]
\centering
\resizebox{\linewidth}{!}{%
\begin{tabular}{lllll}
\toprule
Topology & $\EH$ & $\AX$ & $\VT$ & $\MX$ \\
\midrule
$\Sthree$   & \code{L_ADMISSIBLE} & \code{L_ADMISSIBLE} & \code{L_ADMISSIBLE} & \code{L_CONDITIONALLY_ADMISSIBLE} \\
$\Tthree$   & \code{L_ADMISSIBLE} & \code{L_ADMISSIBLE} & \code{L_ADMISSIBLE} & \code{L_CONDITIONALLY_ADMISSIBLE} \\
$\Nilthree$ & \code{L_ADMISSIBLE} & \code{L_ADMISSIBLE} & \code{L_ADMISSIBLE} & \code{L_CONDITIONALLY_ADMISSIBLE} \\
$\Solthree$ & \code{L_ADMISSIBLE} & \code{L_ADMISSIBLE} & \code{L_ADMISSIBLE} & \code{L_CONDITIONALLY_ADMISSIBLE} \\
\bottomrule
\end{tabular}%
}
\caption{Four-topology admissibility classification.}
\label{tab:four_topology_admissibility}
\end{table}

The label \code{L_ADMISSIBLE} indicates that the Hamiltonian constraint
closes (Section~\ref{sec:hamiltonian:constraint}), the active-torsion
auxiliary closes (\code{AUX_PASS}), and the diagnostic checks
$\Pform=0$ and $\Pint=0$ are passed.  $\EH$ is a pure-curvature reference
to which the P-channels are not applied.  The label
\code{L_CONDITIONALLY_ADMISSIBLE} indicates that the $\MX$ branch becomes
admissible only under the real-auxiliary-branch condition
$\kap^{4}\thetaNY^{2}+1>0$ and the auxiliary-shell interpretation
$\eta=V=0$.  The off-shell $\PintMX$ and $\QNY$ are recorded as
diagnostic / boundary channels but do not obstruct admissibility.

The structural points of this table can be summarised in three statements.
First, the $\EH$ / $\AX$ / $\VT$ / $\MX$ mode pattern agrees across the
four topologies.  Second, the conditions of the conditional admissibility
of $\MX$ are also common to the four topologies (see the auxiliary
determinant in Section~\ref{sec:hamiltonian:constraint}), so topology does
not change the presence or absence of these conditions.  Third, although
$\Solthree$ carries a caveat regarding the compact quotient
(see Appendix~\ref{app:A}), in the main classification table its row
structure coincides with that of the other three topologies.

This topology-robust pattern is the background for the $\chi$-universality
of Section~\ref{sec:bridge}.  Since admissibility itself does not vary with
topology, topology dependence appears not in the admissibility label but
in the diagnostic / orbit quantities $\Ctop$ and $\chi$.  In this sense,
the bridge result of Section~\ref{sec:bridge} is a natural consequence
that rides on top of the admissibility classification.
